% ============================================================
%  CHAPTER 8: TESTING AND EXPERIMENTAL METHODOLOGY
% ============================================================

\chapter{TESTING AND EXPERIMENTAL METHODOLOGY}
\label{chap:testing}

\section{Verification and Validation Strategy}
\label{sec:testing_strategy}

A structured multi-tiered verification and validation (V\&V) strategy was executed to ensure that QuantumSafe VPN satisfies all functional specifications, cryptographic invariants, and mobile OS constraints across five discrete tiers:
\begin{enumerate}
    \item Tier 1: Cryptographic Unit Testing: Mathematical verification of ML-KEM-1024, ML-DSA-65, Keccak SHA3-256, and AES-256-GCM against NIST known-answer test vectors.
    \item Tier 2: Wire Protocol Conformance: Cross-platform binary serialization and shared secret decapsulation verification between Android Kotlin and Linux Python runtimes.
    \item Tier 3: Adversary Injection Testing: Active simulation of Man-in-the-Middle (MitM) spoofing, ciphertext bit-flipping, and packet replay injection.
    \item Tier 4: System Integration Verification: End-to-end evaluation of \texttt{VpnService} packet capture, kernel socket protection via \texttt{protect()}, IP pool leases, and Linux NAT masquerading.
    \item Tier 5: Operational Usability and Robustness: UI state synchronization, background service persistence, abrupt network interface handovers, and DNS leak prevention.
\end{enumerate}

\section{Experimental Testbed Environment}
\label{sec:testbed_environment}

Testing was conducted across a dedicated physical hardware testbed:
\begin{itemize}
    \item Physical Mobile Client: Google Pixel running stock Android 14 (API Level 34, non-rooted) connected via 5\,GHz Wi-Fi (\texttt{11.12.2.82/24}); assigned VPN virtual IP \texttt{10.8.0.2}.
    \item Linux VPN Gateway: Dedicated x86\_64 host (Intel Core i7, 32\,GB RAM) running Kali Linux 2024.2 (Kernel 6.6.9) with Gigabit Ethernet (\texttt{11.12.5.8/24}); virtual interface \texttt{tun0} (\texttt{10.8.0.1/24}, MTU 1400) and \texttt{iptables} NAT masquerading.
    \item Instrumentation Node: Dual-homed promiscuous packet sniffing running Wireshark 4.2.2 and \texttt{tcpdump} tapping both external Wi-Fi transport (TCP 7000) and internal server \texttt{tun0}.
\end{itemize}

\section{Unit and Interoperability Verification}
\label{sec:unit_interoperability_testing}

Six automated Python test suites in \texttt{tests/} execute in 4.18 seconds with a 100\% pass rate, confirming mathematical correctness of lattice key generation, signing, dynamic IP allocation, and tampering detection.

Table~\ref{tab:wire_interop_matrix} summarizes cross-platform wire protocol interoperability between the Android Kotlin client and the Linux Python server.

\begin{table}[htbp]
  \centering
  \small
  \setlength{\tabcolsep}{3pt}
  \caption{Wire Protocol Interoperability and Alignment Verification Matrix.}
  \label{tab:wire_interop_matrix}
  \begin{tabularx}{\linewidth}{@{}p{2.6cm}YYp{2.2cm}@{}}
    \toprule
    \textbf{Protocol Stage} & \textbf{Linux Server (Python)} & \textbf{Android Client (Kotlin)} & \textbf{Status} \\
    \midrule
    Length Framing & 4-byte big-endian (\texttt{!I}) & 4-byte big-endian (\texttt{putInt}) & 100\% Identical \\
    Stage 1: Server Hello & Sends 1,952B ML-DSA-65 PK & Receives 1,952B; pins hash & 100\% Identical \\
    Stage 2: Client Hello & Receives 1,568B ML-KEM PK & Generates \& sends 1,568B PK & 100\% Identical \\
    Stage 3A: Ciphertext & Encapsulates; sends 1,568B $c$ & Decapsulates matching 32B $SS$ & 100\% Identical \\
    Stage 3B: Signature & Signs transcript (3,309B $\sigma$) & Verifies $\sigma$ over transcript & 100\% Identical \\
    Stage 3C: IP Config & Sends \texttt{\{"assigned\_}\allowbreak\texttt{ip":}\allowbreak\texttt{"10.8.0.2"\}} & Parses JSON; configs TUN & 100\% Identical \\
    Stage 4: Key KDF & HKDF-SHA256 ($SS$, salt=$0^{32}$) & HKDF-SHA256 ($SS$, salt=$0^{32}$) & 100\% Identical \\
    Data Framing & \texttt{[4B len] + [12B IV] + Payload + [16B tag]} & \texttt{[4B len] + [12B IV] + Payload + [16B tag]} & 100\% Identical \\
    \bottomrule
  \end{tabularx}
\end{table}

\section{Security Adversary Injection and Leak Testing}
\label{sec:security_adversary_testing}

Active security testing validated protocol resilience against simulated cyber adversaries:
\begin{enumerate}
    \item Ciphertext Bit-Flipping (\texttt{test\_vpn\_tampering.py}): An adversary proxy flipped random bits in the 12-byte IV, payload, and 16-byte Galois MAC tag across 1,000 trials. In 100\% of cases, the receiving endpoint raised an \texttt{AEADBad}\allowbreak\texttt{TagException} or \texttt{InvalidTag} exception, instantly dropping the corrupted frame before decapsulation.
    \item Anti-Replay Counter Injection: Capturing 100 valid encrypted frames ($S \in [50, 150]$) and re-injecting them when the active counter reached $S = 450$ resulted in all 100 packets being dropped ($S_{\text{incoming}} \le S_{\text{last\_received}}$) with zero state disruption.
    \item MitM Rogue Server Substitution: Substituting the server's ML-DSA public key with an attacker keypair caused an immediate client-side fingerprint mismatch in \texttt{SecurityManager}, aborting the connection prior to sending any client key material.
    \item DNS Leak and Routing Audit: Continuous packet sniffing during full-tunnel operation confirmed that 100\% of IPv4 traffic was routed through \texttt{tun0}, with zero cleartext DNS datagrams (UDP port 53) leaking over the physical Wi-Fi interface.
\end{enumerate}

\section{Consolidated System Test Execution Matrix}
\label{sec:test_execution_matrix_sec}

Table~\ref{tab:test_execution_matrix} documents the comprehensive test execution catalogue across all functional, cryptographic, network, and UI test cases.

\begin{table}[htbp]
  \centering
  \scriptsize
  \setlength{\tabcolsep}{2.5pt}
  \caption{Comprehensive System Test Execution Matrix.}
  \label{tab:test_execution_matrix}
  \begin{tabularx}{\linewidth}{@{}p{1.1cm}p{1.6cm}p{3.8cm}Yp{1.1cm}@{}}
    \toprule
    \textbf{ID} & \textbf{Domain} & \textbf{Test Objective} & \textbf{Observed Verification Outcome} & \textbf{Status} \\
    \midrule
    TC-01 & Permissions & \texttt{VpnService} user consent & System dialog displayed; consent granted cleanly & PASS \\
    TC-02 & Socket & Kernel socket loop bypass & \texttt{protect(socket)} returns true; bypasses TUN & PASS \\
    TC-03 & Crypto & ML-KEM-1024 keypair gen & Produced valid 1,568B PK and 3,168B SK & PASS \\
    TC-04 & Crypto & ML-KEM-1024 secret parity & Identical 32B shared secret recovered across nodes & PASS \\
    TC-05 & Crypto & ML-DSA-65 identity signature & 3,309B signature verified over session transcript & PASS \\
    TC-06 & Crypto & Keccak SHA3-256 parity & Bit-for-bit identical digests in Kotlin and Python & PASS \\
    TC-07 & Crypto & AES-256-GCM AEAD enc/dec & 10,000 randomized packets decrypted with valid tags & PASS \\
    TC-08 & Security & Ciphertext bit-flip detection & \texttt{AEADBad}\allowbreak\texttt{TagException}; altered frame dropped & PASS \\
    TC-09 & Security & Anti-replay enforcement & Replayed packets ($S \le S_{\text{last}}$) discarded & PASS \\
    TC-10 & Security & Server identity pinning check & Rogue key rejected; handshake aborted immediately & PASS \\
    TC-11 & Network & Dynamic IP pool lease & Assigned \texttt{10.8.0.2}; subnet \texttt{255.255.255.0} & PASS \\
    TC-12 & Network & Linux NAT packet forwarding & Pings to \texttt{1.1.1.1} forwarded and returned cleanly & PASS \\
    TC-13 & Network & MTU boundary (1400B) & 1400B payload forwarded without fragmentation & PASS \\
    TC-14 & Network & DNS leak prevention & Zero port 53 packets on wlan0; 100\% tunneled & PASS \\
    TC-15 & UI & State badge transition & DISCONNECTED $\to$ CONNECTING $\to$ CONNECTED & PASS \\
    TC-16 & UI & Telemetry throughput & Uplink/downlink rates update smoothly in UI & PASS \\
    TC-17 & Storage & Keystore configuration & Endpoint settings survive reboot via AES-256 & PASS \\
    TC-18 & Lifecycle & Graceful disconnection & TUN descriptor closed; socket cleanly terminated & PASS \\
    \bottomrule
  \end{tabularx}
\end{table}

All 18 executed test cases achieved a 100\% pass rate, confirming the functional validity, cryptographic integrity, and operational robustness of QuantumSafe VPN.
